\chapter{Instrukcja uruchomienia projektu}
\label{cha:uruchomienie}

\section{Wymagania systemowe}
\label{sec:wymaganiaUruch}

\begin{itemize}
  \item \textbf{System operacyjny:} Windows 10 lub Windows 11.
  \item \textbf{Środowisko uruchomieniowe:} platforma .NET (wersja zgodna z projektem) ze wsparciem WPF.
  \item \textbf{Serwer bazy danych:} PostgreSQL (zalecana wersja 13 lub nowsza).
  \item \textbf{Dodatkowe narzędzia:} środowisko programistyczne wspierające .NET (np. Visual Studio)
        oraz narzędzie do zarządzania bazą danych (np. pgAdmin).
\end{itemize}

\section{Konfiguracja bazy danych}
\label{sec:konfBaza}

Aplikacja łączy się z bazą danych PostgreSQL przy użyciu parametrów zdefiniowanych w klasie
\texttt{AppDbContext}: serwer \texttt{localhost}, port \texttt{5432}, baza \texttt{testownik},
użytkownik \texttt{postgres}. Przed pierwszym uruchomieniem należy zapewnić działający serwer
PostgreSQL z dostępem do tych parametrów (w razie potrzeby parametry połączenia można dostosować
w kodzie klasy \texttt{AppDbContext}).

Ręczne tworzenie bazy ani tabel nie jest wymagane - schemat zakładany jest automatycznie. Przy
starcie aplikacja wywołuje metodę \texttt{Database.Migrate()}, która tworzy bazę i nakłada wszystkie
migracje Entity Framework Core, a następnie klasa \texttt{DbSeeder} wypełnia ją danymi początkowymi
(przykładowy nauczyciel, uczniowie, grupy oraz testy).

\section{Konfiguracja środowiska i uruchomienie aplikacji}
\label{sec:uruchAplikacji}

Kolejne kroki uruchomienia:

\begin{enumerate}
  \item uruchomić serwer PostgreSQL i upewnić się, że jest dostępny na \texttt{localhost:5432};
  \item otworzyć projekt w środowisku Visual Studio (plik rozwiązania \texttt{ProjektPO2.slnx})
        i uruchomić aplikację (lub wykonać polecenie \texttt{dotnet run} w katalogu projektu);
  \item przy pierwszym uruchomieniu baza zostanie utworzona i zasilona danymi demonstracyjnymi;
  \item zalogować się przy użyciu jednego z kont demonstracyjnych.
\end{enumerate}

\section{Dane logowania testowego}
\label{sec:daneLogowania}

Po zasileniu bazy danymi początkowymi dostępne są konta demonstracyjne (hasło dla wszystkich:
\texttt{demo1234}):

\begin{itemize}
  \item nauczyciel - login \texttt{m.dabrowski};
  \item uczniowie - loginy \texttt{anna.k}, \texttt{piotr.n}, \texttt{maria.w} oraz kolejni.
\end{itemize}
